\section{Dynamics}
\label{sec:dynamics}

In this section, we look at a simple dynamic that leads to approximate RARE prices.

\subsection{Limiting RARE}
\label{sec:lim-rare}

We begin by giving an alternative formulation of RARE,
which would be more useful for analyzing dynamics.
To do so, we define each seller $j$'s \emph{optimal deviation revenue} $V_j(\pvec_{-j})$.

\begin{definition}[Optimal deviation revenue]
\label{defn:vj}
Let $\Ical = ([n], [m], (u_i)_{i=1}^n, (b_i)_{i=1}^n)$ be a market,
$A: \mathbb{R}_{\ge 0}^m \to [0, 1]^{n \times m}$ be an allocation rule,
and $\pvec \in \mathbb{R}_{\ge 0}^m$ be a price vector.
When seller $j$ deviates to price $q_j \ge 0$,
define the other sellers' set of individually rational (IR) responses as
\[ \IR^A_j(q_j, \pvec_{-j}) \defeq \{\qvec_{-j} \in \mathbb{R}_{\ge 0}^{[m] \setminus \{j\}}:
    R^A_k(\qvec) \ge R^A_k((q_j, \pvec_{-j}))\; \forall k \in [m] \setminus \{j\}\}. \]
Note that this set is non-empty, since $\pvec_{-j} \in \IR^A_j(q_j, \pvec_{-j})$.
%
Define the deviator $j$'s \emph{worst-case revenue} at deviation $q_j$ as
\[ D^A_j(q_j, \pvec_{-j}) \defeq \inf_{\qvec_{-j} \in \IR^A_j(q_j, \pvec_{-j})} R^A_j(\qvec), \]
and her optimal deviation revenue as
\[ V^A_j(\pvec_{-j}) \defeq \sup_{q_j \in \mathbb{R}_{\ge 0}} D^A_j(q_j, \pvec_{-j}). \]
When $A$ is clear from context, we simply write $\IR_j$, $D_j$, and $V_j$
instead of $\IR^A_j$, $D^A_j$, and $V^A_j$.
\end{definition}

Intuitively, it might seem that a price vector $\pvec$ is RARE iff
$R_j(\pvec) \ge V_j(\pvec_{-j})$ for every seller $j \in [m]$,
giving us an alternative formulation of RARE.
But this formulation is \emph{not} definitionally equivalent to RARE, for a subtle reason.
In particular, for every deviation $q_j$ by seller $j$,
RARE requires an IR response $\qvec_{-j}$
such that the new revenue $R_j(\qvec)$ is at most $R_j(\pvec)$.
Our new definition, instead, requires a \emph{sequence of} IR responses for which $j$'s new revenue
converges to at most $R_j(\pvec)$.
We call this new definition \emph{limiting-RARE}.

\begin{definition}[limiting-RARE]
\label{defn:lim-RARE}
Consider a market $\Ical$ and allocation rule $A$.
Then for any $\delta \ge 0$ and $\kappa \ge 1$, the price vector $\pvec$
is called $\delta$-additive $\kappa$-\limRARE{} if for every seller $j \in [m]$, we have
\[ R^A_j(\pvec) \ge \frac{V^A_j(\pvec_{-j}) - \delta}{\kappa}. \]
\end{definition}

The following result shows that RARE and \limRARE{} are the same up to
an infinitesimally small additive constant.
We defer the proof to \cref{sec:lim-rare-equiv}.

\begin{restatable}{lemma}{thmRareEquivDefn}
\label{thm:rare-equiv-defn}
For any $\delta \ge 0$ and $\kappa \ge 1$,
a price vector $\pvec$ is $\delta$-additive $\kappa$-\limRARE{} iff
for every $\eps > 0$, it is $(\delta + \eps)$-additive $\kappa$-RARE.
Moreover, a $\delta$-additive $\kappa$-RARE price vector is also
$\delta$-additive $\kappa$-\limRARE.
\end{restatable}

\subsection{The Dynamic}
\label{sec:dynamics:defn}

Fix a continuous allocation rule $A$. The sellers initially have prices $\pvec^{(0)} = \vecZero$.
At each time step $t \in \N$, a seller $j$ is called \emph{unsatisfied} if
$V^A_j(\pvec^{(t-1)}_{-j}) > \kappa\nu p^{(t-1)}_j$.
Denote the set of all unsatisfied sellers by $J^{(t)}$.
If $J^{(t)} = \emptyset$, i.e., there are no unsatisfied sellers,
then let $\pvec^{(t)} \defeq \pvec^{(t-1)}$.
Otherwise, pick an arbitrary non-empty subset $S^{(t)} \subseteq J^{(t)}$, and update the prices as
\[ p^{(t)}_j \defeq \begin{cases}
\displaystyle \frac{V^A_j(\pvec^{(t-1)}_{-j})}{\kappa\nu} & \text{ if } j \in S^{(t)}
\\[0.8em] p^{(t-1)}_j & \text{ if } j \notin S^{(t)}
\end{cases}. \]

We will soon prove that $R_j(\pvec^{(t)}) = \nu p^{(t)}_j$ for all $t$ and $j$, which means that when $\kappa = 1$,
each seller either retains her old price, or increases her price to get her optimal deviation revenue,
making this a very natural dynamic.

Note that if $J^{(t)} = \emptyset$, then the prices freeze permanently.
Otherwise, prices of goods $S^{(t)}$ strictly increased, and the other goods had the same price as the previous round.
By picking $S^{(t)}$ appropriately, we can model different types of dynamics:
if $S^{(t)} = J^{(t)}$ for all $t$, then sellers update prices synchronously,
and if $|S^{(t)}| = 1$ for all $t$, then sellers update prices one-at-a-time.

\subsection{Analyzing the Dynamic}

The $\kappa$-RARE price vector whose existence we established in
\cref{thm:rare-exists-continuous} plays a key role in the analysis.
We begin by proving a useful property of that price vector:
it helps us upper bound the optimal deviation revenue of the price vectors it dominates.

\begin{lemma}
\label{thm:vbound-pstar}
There exists $\pvec^* \in \beta\Delta_m^{\circ}$ such that for every $\pvec \in \prod_{k=1}^m [0, p^*_k]$,
and seller $j \in [m]$, we have $V^A_j(\pvec_{-j}) \le \kappa\nu p^*_j$.
\end{lemma}
\begin{proof}
Let $\pvec^*$ be the $\kappa$-RARE price vector from \cref{thm:rare-exists-continuous}.
Pick any seller $j$ and deviation $q_j \ge 0$.
By \cref{thm:rare-exists-continuous}, there exists a response $\qvec_{-j}$ such that
\begin{tightenum}
\item $R_j(\qvec) \le \kappa R_j(\pvec^*) = \kappa\nu p^*_j$.
\item For each responder $k \in [m] \setminus \{j\}$, we have $R_k(\qvec) \ge \nu p_k^*$.
\end{tightenum}

For any responder $k \in [m] \setminus \{j\}$, we have
$R_k((q_j, \pvec_{-j})) \le \nu p_k \le \nu p_k^* \le R_k(\qvec)$.
Thus, $\qvec_{-j} \in \IR_j(q_j, \pvec_{-j})$, and so,
$D_j(q_j, \pvec_{-j}) \le R_j(\qvec) \le \kappa \nu p^*_j$.
Hence, $V_j(\pvec_{-j}) \le \kappa \nu p^*_j$.
\end{proof}

We now show that in each iteration of the dynamic, the prices are upper bounded by $\pvec^*$.

\begin{lemma}
Let $(\pvec^{(t)})_{t \in \N \cup \{0\}}$ be the dynamic defined in \cref{sec:dynamics:defn}.
Let $\pvec^*$ be as defined in \cref{thm:vbound-pstar}.
Then for all $t \in \N \cup \{0\}$ and $j \in [m]$, we have
$p^{(t)}_j \le p^{(t+1)}_j \le p^*_j$ and $R_j(\pvec^{(t)}) = \nu p^{(t)}_j$.
\end{lemma}
\begin{proof}
The monotonicity of $\pvec^{(t)}$ in $t$ follows easily from its definition,
i.e., $p^{(t)}_j \le p^{(t+1)}_j$ for all $j \in [m]$.

For all $t \in \N$ and $j \in [m]$, using \cref{thm:vbound-pstar}, we get
\[ p^{(t)}_j \le \max\bigg(p^{(t-1)}_j, \frac{V_j(\pvec^{(t-1)}_{-j})}{\kappa\nu}\bigg)
    \le \max(p^{(t-1)}_j, p^*_j). \]
Then we can prove $p^{(t)}_j \le p^*_j$ for all $j \in [m]$ and $t \in \N \cup \{0\}$ by induction on $t$,
since $\pvec^{(0)} = \vecZero$.

Moreover,
\[ \sum_{j=1}^m p^{(t)}_j \le \sum_{j=1}^m p^*_j = \beta, \]
so $R_j(\pvec^{(t)}) = \nu p^{(t)}_j$ for all $j$ by \cref{thm:accounting}.
\end{proof}

Observe that if the dynamic reaches equilibrium (i.e., the prices stop changing),
then the prices are $\kappa$-\limRARE{}.

\begin{observation}
Let $\pvec^{(0)} = \vecZero$, and for $t \in \N$, let $p^{(t)}$, $J^{(t)}$, and $S^{(t)}$
be as defined in \cref{sec:dynamics:defn}.
Suppose the prices do not change in some iteration $T+1$, i.e., $\pvec^{(T+1)} = \pvec^{(T)}$.
Then $R_j(\pvec^{(T)}) = \nu p^{(T)}_j \ge V_j(\pvec^{(T)}_{-j})/\kappa$ for all $j \in [m]$,
i.e., $\pvec^{(T)}$ is $\kappa$-\limRARE{}.
\end{observation}

Note that the prices $\pvec^{(T)}$ above need not be the same as $\pvec^*$.
There may be multiple $\kappa$-\limRARE{} price vectors.

Even if the dynamic does not reach equilibrium, we show that after sufficiently many rounds,
we get $\kappa$-RARE with a small additive approximation
if $S^{(t)}$ always includes the \emph{most unsatisfied} seller.

\begin{theorem}
\label{thm:dynamics-convergence}
Let $\pvec^{(0)} = \vecZero$, and for $t \in \N$, let $p^{(t)}$, $J^{(t)}$, and $S^{(t)}$
be as defined in \cref{sec:dynamics:defn}.
In every iteration $t \in \N$ where $J^{(t)} \neq \emptyset$, suppose $S^{(t)}$ always includes
some seller $j$ maximizing $V_j(\pvec^{(t-1)}_{-j})/(\kappa\nu) - p^{(t-1)}_j$.
Then for any $\delta > 0$:
\begin{enumerate}
\item (an approximate RARE occurs soon) $\exists t \le \floor{B/\delta}$ such that
    $\pvec^{(t)}$ is a $\delta$-additive $\kappa$-RARE.
\item (eventually prices remain approximately RARE) $\exists T \in \N \cup \{0\}$ such that
    for all $t > T$, $\pvec^{(t)}$ is a $\delta$-additive $\kappa$-RARE.
\end{enumerate}
\end{theorem}
\begin{proof}
For any $t \in \N \cup \{0\}$, let
\[ g_t \defeq \max_{j \in [m]} \max(0, V_j(\pvec^{(t)}_{-j}) - \kappa R_j(\pvec^{(t)})). \]
Then $\pvec^{(t)}$ is $g_t$-additive $\kappa$-\limRARE.
%
Since $R_j(\pvec^{(t)}) = \nu p^{(t)}_j$ for all $t$ and $j$,
we get $g_t = 0$ if $J^{(t+1)} = \emptyset$. Otherwise,
\begin{align*}
\kappa\nu \max_{j \in [m]} (p^{(t+1)}_j - p^{(t)}_j)
   &= \kappa\nu \max_{j \in S^{(t+1)}} (p^{(t+1)}_j - p^{(t)}_j)
    = \max_{j \in S^{(t+1)}} V_j(\pvec^{(t)}_{-j}) - \kappa R_j(\pvec^{(t)})
\\ &= \max_{j \in [m]} \max(0, V_j(\pvec^{(t)}_{-j}) - \kappa R_j(\pvec^{(t)}))
    = g_t.
\end{align*}
(Here the last equality follows from $S^{(t+1)}$ containing the most unsatisfied seller.)

Now for any $T \in \N \cup \{0\}$, sum this equality over all $t \in [T] \cup \{0\}$ to get
\[ \sum_{t=0}^T g_t = \kappa\nu \sum_{t=0}^T \max_{j \in [m]} (p^{(t+1)}_j - p^{(t)}_j)
    \le \kappa\nu \sum_{j=1}^m \sum_{t=0}^T (p^{(t+1)}_j - p^{(t)}_j)
    = \kappa\nu \sum_{j=1}^m p^{(T+1)}_j \le \kappa\nu\beta = B. \]
Pick any $\delta > 0$.
\begin{enumerate}
\item Let $T \defeq \floor{B/\delta}$. Then for some $t \in [T] \cup \{0\}$, we have
    $g_t \le B/(T+1) < \delta$, which implies that $\pvec^{(t)}$ is a $\delta$-additive $\kappa$-RARE.
    (Note that we get an approximate RARE, not just an approximate \limRARE{},
    since $g_t$ is \emph{strictly} less than $\delta$.)
\item There are at most $\floor{B/\delta}$ values of $t$ for which $g_t \ge \delta$.
    Let $T$ be the largest such that $g_T \ge \delta$ (and let $T$ be 0 if $g_t < \delta$ for all $t$).
    Then for all $t > T$, we have $g_t < \delta$, i.e., $\pvec^{(t)}$ is a $\delta$-additive $\kappa$-RARE.
    \qedhere
\end{enumerate}
\end{proof}

Intuitively, if we don't require $S^{(t)}$ to always include the most unsatisfied seller,
one seller might get stuck at a particularly bad price while the other sellers keep improving.
